\subsection{Closing the full rank-one $WT^2$ family}\label{sec:rankoneformal}

\begin{proposition}\label{prop:rankoneformal}
For $f\ne0$ and $\deg G\leq3$, the polynomial
\[
 P=W(a+bZ+cT+fT^2)+G(Z,T)
\]
is pure if and only if $b=0$ and $G=0$. Its pure members obey
$m>\deg Q$.
\end{proposition}
\begin{proof}
Put $q=1-bt$ and $A(T)=a+cT+fT^2$. Exact circular contraction
gives
\[
 J(p)=\frac1q\E_T\!\left[p(T)
             \exp\{tG(tA/q,T)\}\right],\qquad J(1)=\Phi_1.
\]
Gaussian integration by parts says that $J$ annihilates
$\partial_T+\varphi_T$; direct parameter differentiation makes it
horizontal for the polynomial phase and normalization
\[
 \varphi=tG(tA/q,T)-T^2/2,\qquad h=1/q.
\]
The phase has degree at least two, since it equals $-T^2/2$ at $t=0$.

Suppose first $b=0$ and write $G=\sum_{j\geq0}Z^jG_j(T)$.
If a term $A^jG_j$ attains the actual degree $D\geq2$ of the phase,
the leading coefficient is a polynomial in $t$ containing a
nonzero positive power. Distinct $j$ give distinct powers $t^{j+1}$,
so they cannot cancel identically. Its order at infinity is negative,
whereas $v_\infty(h)=0$. Criterion~\eqref{eq:leadingcriterion}
excludes purity. Because $\deg A=2$, every nonzero $G_j$ with
$j\geq1$ attains some degree at least two. The only remaining
possibility is $G=\alpha T+\beta$. Its exact unit series is
$\exp(\beta t+\alpha^2t^2/2)$, which is one only if $G=0$.

Suppose now $b\ne0$. Divide $G$ by $S=bZ+A(T)$ as a polynomial
in $T$:
\[
 G=R_0(Z)+TR_1(Z)+S(\alpha Z+\beta T+\gamma).
\]
The remainder has basis $1,T,Z,ZT,Z^2,Z^2T,Z^3$. At $Z=-A/b$
these monomials have distinct leading degrees $0,1,2,3,4,5,6$.
Thus if $V(T)=R(-A/b,T)$ has degree $n$, then
$\deg R(sA,T)\leq n$ for every scalar $s$; a constant $V$
forces $R$ to be that same constant. Write $g=tA/q$. The phase is
exactly
\[
 tR(g,T)+\alpha g^2+\beta gT+\gamma g-T^2/2.
\]
Since $g=-A/b+O(t^{-1})$, its first term is $tV+O(1)$, with
all these remainders of degree at most $n$.

If $\alpha\ne0$, the limiting remaining phase has degree four
with leading coefficient $\alpha f^2/b^2$. If $\alpha=0$,
$\beta\ne0$, it has degree three with leading coefficient
$-\beta f/b$. Otherwise its quadratic coefficient tends to
$-\gamma f/b-1/2$. A leading term from $tV$ has valuation $-1$;
a higher-degree nonzero limiting term has valuation zero.
The only flattened case has
$\alpha=\beta=0$, $\gamma=-b/(2f)$. Its exact quadratic
coefficient outside $tR$ is $-1/(2q)$, of valuation one.
Thus in every case the actual leading coefficient $a_D$ has
$v_\infty(a_D)\leq1$. This also covers $V$ constant or linear.
But $v_\infty(h)=1$ and $D\geq2$, so
$Dv_\infty(h)-v_\infty(a_D)\geq D-1>0$.
Purity is impossible by the same criterion.

Conversely $W A(T)$ has strictly negative circular weight.
For a monomial $Z^iW^jT^k$, a mixed moment can be nonzero only
when the number of $W$ factors balances the number of $Z$ factors;
in particular it vanishes for $m>\deg Q$.
\end{proof}
